% !TeX root = ../深度学习与大语言模型技术原理.tex
% 由主文档通过 \input 载入；请编译主文档。
\setcounter{chapter}{17} % 保留原书章号。
\chapter{智能体强化学习、代码与环境模型}
\label{ch:agentlearning}

\section{部分可观察环境中的策略学习}

工具智能体的真实环境状态可能包含文件、网页、数据库或隐含任务条件，而模型只看见部分观察。因此常用部分可观察马尔可夫决策过程（partially observable Markov decision process, POMDP）建模。策略可依赖历史 $h_t=(o_1,a_1,\ldots,o_t)$ 或它的记忆表示；如果摘要丢失关键状态，就不能默认摘要仍满足马尔可夫性质。

智能体强化学习（agent RL）优化多次动作后的任务回报，工具使用强化学习（tool-use RL）强调选择工具与参数的学习。轨迹概率包含策略动作与环境转移；通常只对模型生成动作的对数概率求策略梯度，工具返回文本不是模型动作，不能一并当作需要提高概率的目标。把全部日志直接做语言建模，可能让模型学会伪造观察。

训练环境要支持可重置初始状态、隔离副作用和稳定评分。异步 rollout 提高利用率，却可能使用过时策略；需要记录采样版本、控制滞后并说明校正方法。长度差异大的轨迹还会造成负载不均，训练吞吐不能只按任务数统计。

\section{稀疏奖励与势函数塑形}

若只有任务最终成功才得到奖励，长轨迹中的有用动作很难归因。中间奖励可以来自有效证据、通过子测试或达到可验证里程碑，但奖励设计不当可能鼓励无意义动作。一个有理论依据的形式是势函数塑形（potential-based reward shaping）：
\[
r'_t=r_t+\gamma\Phi(s_{t+1})-\Phi(s_t).
\]
对有限时域 $t=0,\ldots,T-1$，塑形项折扣和为 $-\Phi(s_0)+\gamma^T\Phi(s_T)$。若终止势为零且初始状态固定，它只是与策略无关的常数，故不改变期望回报的最优策略。若终止势、时域或部分观察处理不满足条件，就不能直接援引这个结论。

实际中可用最终环境成功验证器判断任务完成，再把过程指标作为诊断或谨慎的辅助奖励。任务奖励上升后，还应检查动作数量、工具滥用、环境损坏和未覆盖约束，避免模型学会通过评分漏洞而非解决问题。

\section{代码模型与软件工程智能体}

代码语言模型可以用代码预训练、指令微调和执行反馈学习。中间填充（fill-in-the-middle, FIM）通过特殊标记重排前缀、后缀和中间缺口，训练根据两侧上下文补全内容；它依然可以用自回归损失实现。仓库级任务还需要理解符号定义、依赖、调用链、测试和构建环境，而非只处理孤立函数。

软件工程智能体围绕问题描述检索代码、读取文件、提出补丁、运行相关测试并依据失败反馈修正。仓库上下文应结合精确符号搜索与语义检索；错误地把相似文件当成实际入口会产生看似合理却不起作用的补丁。补丁应作为可审查产物，记录改动范围和测试证据。

单元测试奖励只覆盖被执行的输入与性质。测试通过不证明程序完全正确，更不能允许模型修改测试来伪造成功。训练和评测应隔离隐藏测试，固定环境与依赖，区分编译失败、运行超时、断言失败和环境故障。代码执行必须位于合适的沙箱与权限范围内。

\section{计算机操作与视觉定位}

计算机操作智能体（computer-use agent）根据屏幕图像、界面树或状态信息选择鼠标、键盘和应用动作。视觉定位（visual grounding）把语言目标关联到实际控件或区域。分辨率变化、滚动偏移、动态布局和遮挡都会影响坐标，旧截图上的位置不能无条件用于新状态。

动作前应确认关键界面条件，动作后检查预期状态变化。环境非确定性包括网络延迟、弹窗、异步加载和焦点变化，恢复策略应区别“动作未执行”与“执行后反馈延迟”。表单提交等副作用动作尤其需要可观测确认，而不能仅凭点击坐标推断成功。

评测可使用受控网页或桌面任务，记录成功率、有效动作数、误操作与恢复能力。多模态模型改善感知接口，但不能替代状态管理和权限控制；视觉能力与任务规划应分开诊断。

\section{世界模型与环境模型}

世界模型（world model）或环境模型（environment model）学习状态转移与反馈，例如 $p_\phi(o_{t+1},r_t\mid h_t,a_t)$，以支持预测、模拟和规划。它可以在真实执行前估计动作后果，或用模拟轨迹补充训练，但模型中的成功并不等于真实环境成功。

若一步预测有误差，长时域滚动可能不断偏离真实可达状态；规划还可能主动选择模型最乐观却最不准确的区域。可通过短期规划、模型集合、不确定性估计与真实反馈校正缓解，但必须明确哪些结果来自模拟、哪些经过真实验证。

在 LLM 场景中，用语言描述“下一步可能发生什么”只是环境模型的一种表达形式。是否能作为可靠规划器，取决于预测校准、动作条件覆盖和闭环验证，而不是生成描述是否连贯。

